%% kapa: exam=1 round=31 year=2020 period=1 form=A subject=real_estate subject_ko=부동산학원론 seq=1 number=81
%% source: Q-Net 감정평가사 기출문제 (cst00309) artlSeq=5207219 file=extracted/1008/5207219/A형/A형/2020년도 제31회 감정평가사 1차 시험_1교시_A형.pdf
%% conversion: kordoc 4.24.51 -> tools/gen1.py (auto); review=proofread-sonnet; flags=none
\begin{question}{81}
토지의 특성에 관한 설명으로 옳지 않은 것은?
\begin{choices}
  \item 부동성으로 인해 지역분석을 필요로 하게 된다.
  \item 용도의 다양성은 최유효이용의 판단근거가 된다.
  \item 영속성은 부동산활동에 대해서 장기적 배려를 필연적으로 고려하게 한다.
  \item 합병ㆍ분할의 가능성은 토지의 이행과 전환을 가능하게 한다.
  \item 개별성으로 인해 일물일가의 법칙이 적용되지 않고, 부동산시장에서 부동산상품 간에 완벽한 대체는 불가능하다.
\end{choices}
\end{question}
\answer{4}
